%!TEX root = ../main.tex
\section{Source Design, Ingestion, and Alignment}
\label{sec:acquisition}

Preparation starts before the first byte is recorded.  An acquisition interface
determines what states can be visited, how accurately actions are labeled, which
failures are observable, and how expensive later alignment will be.  This section
therefore includes source design when it intentionally controls compatibility or
coverage, while leaving hardware design itself outside our scope.

\subsection{Sources Have Complementary Training Roles}

Real-robot teleoperation offers the strongest observation--action--physics
alignment but is expensive and biased by the available robot, operator, and task
script.  DROID reduces site bias through a common mobile manipulator and protocol
deployed across 18 institutions, yielding roughly 76,000 demonstrations in
diverse scenes~\cite{khazatsky2024droid_2403_12945}.  BridgeData V2 emphasizes an
open, multi-task collection on a consistent platform~\cite{walke2023bridgedata_2308_12952}.
Open X-Embodiment takes the opposite integration route: convert existing
datasets into a common interface and train across institutions and robots
~\cite{collaboration2023open_2310_08864}.  The former controls acquisition while
the latter accepts heterogeneity and moves work into post-hoc alignment.
RoboTurk decouples demonstrators from the robot through remote crowdsourcing,
making operator latency and identity preparation metadata rather than incidental
noise~\cite{mandlekar2018roboturk_1811_02790}.  AXIS further integrates
community collection, task generation, filtering, smoothing, and augmentation as
a growable engine~\cite{zhao2026axis_2607_21588}; as a recent preprint, it is
evidence of an emerging system pattern rather than a mature standard.

Portable human-operated interfaces move collection away from continuous target
robot access.  UMI reproduces relevant gripper geometry, recovers 6-DoF motion,
and reduces observation mismatch at deployment~\cite{chi2024universal_2402_10329}.
Human videos offer much greater semantic and environmental breadth, but usually
lack robot states, metric scale, actions, forces, and reliable task boundaries;
their use spans representation pretraining, observation transfer, and action
reconstruction~\cite{ma2026robot_2604_27621}.  EgoMimic co-trains aligned
egocentric human and robot demonstrations, whereas HOP reconstructs 3-D
hand--object interaction and retargets it into robot experience
~\cite{kareer2024egomimic_2410_24221,singh2024hand_2409_08273}; both illustrate
that ``human video'' becomes action supervision only through explicit geometric
and embodiment assumptions.  Simulation offers exact state and
controllable long-tail generation but inherits asset, rendering, dynamics, and
contact mismatch.  Autonomous execution, as in AutoRT, produces records from the
current policy in the actual environment, directly exposing its support gaps but
also concentrating data around its existing competence
~\cite{ahn2024autort_2401_12963}.

No source dominates along correctness, diversity, scale, and target alignment.
The more useful abstraction is a \emph{role}: web/human video for broad visual and
task semantics; portable human demonstrations for motion and temporal structure;
multi-robot corpora for shared visuomotor priors; target-robot teleoperation for
control and contact calibration; successful task demonstrations for adaptation;
and failed/recovery execution for post-training and safety.  Routing by role
prevents ``lower-fidelity'' data from being discarded when it supplies a
different signal, and prevents weakly aligned data from silently supervising
precise target actions.

\subsection{Coverage Must Be Designed, Not Counted}

Episode count collapses several axes: task concepts, objects, backgrounds,
viewpoints, initial conditions, operators, embodiments, success/failure modes,
and transition density.  Repeating a narrow configuration can reduce estimator
variance without expanding support; adding a new environment can expand support
but also introduce a confounded action convention.  The acquisition target should
therefore be a structured coverage tensor rather than a global quota.  A practical
engine tracks cells such as $(\text{skill},\text{object},\text{scene},
\text{embodiment},\text{outcome})$, uncertainty in each cell, and marginal cost.
Controlled real-robot scaling experiments similarly find that, after enough
repetition, expanding object and environment diversity is more valuable than
additional demonstrations of the same configuration in their task setting
~\cite{lin2024data_2410_18647}.  The fitted scaling law should not be transferred
unchanged across learners or collection protocols, but its factorized experiment
is the relevant form of evidence.

Active acquisition then chooses a candidate task $q$ by an objective such as
\begin{equation}
q^*=\arg\max_q
\frac{\widehat{\Delta\utility}(q)+\alpha\,\Delta\mathrm{coverage}(q)
-\beta\,\mathrm{risk}(q)}{\mathrm{cost}(q)}.
\end{equation}
Uncertainty alone is unsafe: it can select sensor corruption, infeasible tasks,
or states outside the robot's recovery envelope.  Value must be gated by
executability and safety, and by whether the resulting record can be labeled well
enough to train from.  Autonomous engines make this objective operational, but
most current systems still use hand-designed task proposals and success filters
rather than calibrated marginal utility.

\subsection{Ingestion and Technical Quality Control}

Technical validity is the cheapest part of the evidence ladder and should be
resolved before semantic scoring.  A minimum ingestion contract checks:

\begin{itemize}[leftmargin=*]
    \item file readability, schema/version compatibility, units, shapes, and
    required fields;
    \item monotonic clocks, dropped frames, cross-stream synchronization, and
    command-to-actuation delay;
    \item camera intrinsics/extrinsics, robot description, joint limits, reset
    boundaries, and controller identity;
    \item image blur/exposure/occlusion and sensor saturation, missingness, or
    implausible jumps; and
    \item action range, idle fraction, clipping, jerk, collisions, and emergency
    interventions.
\end{itemize}

These checks should emit repair decisions and uncertainty, not only pass/fail.
Interpolation can repair short sensor gaps but may fabricate contact transitions;
smoothing can remove teleoperation jitter but also erase deliberate high-frequency
corrections.  The unmodified raw stream, repair parameters, and reason code must
remain in lineage.  In contact-rich multimodal data such as RH20T, time alignment
across vision, proprioception, force/torque, audio, and touch is itself a semantic
precondition, not mere storage hygiene~\cite{fang2023rh20t_2307_00595}.

\subsection{Four Layers of Cross-Embodiment Alignment}

Common schemas enable loading and batching, but shape compatibility does not
imply supervision compatibility.  Cross-embodiment alignment has at least four
layers.

\textbf{Structural alignment} maps fields, modalities, missing degrees of freedom,
and gripper or hand state into a common schema.  Fixed action slots plus an
embodiment mask make heterogeneity explicit, but padding only states that a degree
of freedom is absent.

\textbf{Spatial and control alignment} specifies whether actions are joint
positions, velocities, absolute or delta end-effector poses, forces, or discrete
tokens; the reference frame; rotation convention; units; controller gains; and
whether a command is a target or an executed state.  Converting to a camera- or
end-effector-centric representation can remove base-placement variation, but
calibration error then becomes label noise.  Some mappings are many-to-one:
compressing a dexterous hand into ``gripper openness'' loses interaction
semantics and is not invertible.

\textbf{Temporal alignment} reconciles sensor rate, controller rate, observation
latency, action horizon, interpolation, and chunk semantics.  A visually aligned
frame paired with a command issued several control cycles later teaches a
different conditional distribution.  Resampling must be anti-aliased, contact
boundaries should not be averaged away, and both command and measured execution
times should be retained when delays vary.

\textbf{Behavioral alignment} asks whether the same normalized action has the
same consequence under different morphology, kinematics, compliance, and
controllers.  This is the least solved layer.  An embodiment-conditioned action
encoder can incorporate kinematic graphs and control metadata, while cycle or
rollout consistency can test whether mapped actions reproduce the intended state
transition.  Open X-Embodiment demonstrates the benefit of a common training
interface, but its heterogeneity also motivates mixture and learner-aware methods
such as Re-Mix~\cite{collaboration2023open_2310_08864,hejna2024re_2408_14037}.

\subsection{Jointly Transform All Coupled Modalities}

Any spatial or temporal transformation must act consistently on observations,
states, actions, language, and success conditions.  Mirroring an image without
mirroring lateral actions and ``left/right'' language creates confident label
error.  Changing camera pose requires transforming camera-frame actions.  Time
reversal is valid only for a reversible segment and an appropriately reversed
task description; it is generally invalid for pouring, breaking, or frictional
contact.  We call this the \emph{co-transformation principle}: an augmentation is
valid only if it preserves the coupled state transition and task semantics, not
merely image plausibility.

Finally, preparation artifacts need a data contract: source and license,
collection context, robot and controller, coordinate frames, clocks, transformations,
automatic and human labels, confidence, exclusion/routing reasons, train/test
relationships, and version hashes.  Without this contract, later influence
scores, failure reports, and model versions cannot be traced back to the decision
that produced a training example.
